\documentclass[tikz,border=10pt]{standalone}
\usepackage{fontspec}
\usepackage{xeCJK}
\IfFontExistsTF{Noto Sans TC}{\setmainfont{Noto Sans TC}\setCJKmainfont{Noto Sans TC}}{\setmainfont{Microsoft JhengHei}\setCJKmainfont{Microsoft JhengHei}}
\xeCJKsetup{CJKglue={\hskip 0.08em plus 0.02em}, CJKecglue={\hskip 0.11em plus 0.02em}}
\usetikzlibrary{arrows.meta,positioning,calc}
\definecolor{paper}{HTML}{FFFDF8}
\definecolor{navy}{HTML}{111827}
\definecolor{muted}{HTML}{4B5563}
\definecolor{line}{HTML}{CBD5E1}
\definecolor{blue}{HTML}{1D4ED8}
\definecolor{bluebg}{HTML}{DBEAFE}
\definecolor{gold}{HTML}{B45309}
\definecolor{goldbg}{HTML}{FEF3C7}
\definecolor{red}{HTML}{B91C1C}
\definecolor{redbg}{HTML}{FEE2E2}
\definecolor{green}{HTML}{047857}
\definecolor{greenbg}{HTML}{D1FAE5}
\tikzset{
  title/.style={font=\fontsize{15}{20}\selectfont,text=navy,align=center},
  note/.style={font=\fontsize{9.5}{13}\selectfont,text=muted,align=center},
  card/.style={rounded corners=13pt,line width=1pt,align=center,inner xsep=14pt,inner ysep=10pt,font=\fontsize{10.2}{13.8}\selectfont,text=navy}
}
\begin{document}
\begin{tikzpicture}[x=1cm,y=1cm,>=Latex,line cap=round,line join=round]
\path[fill=paper] (0,0) rectangle (16,9);

\node[title,text width=13cm] at (8,8.14) {分數變高，不一定代表事情變好};
\node[note,text width=11.8cm] at (8,7.52) {系統會追逐我們給它的訊號；訊號太窄，捷徑就會像進步。};

\node[card,fill=bluebg,draw=blue,text width=2.45cm] (signal) at (2.35,4.85) {獎勵訊號};
\node[card,fill=goldbg,draw=gold,text width=2.45cm] (opt) at (5.95,4.85) {努力優化};
\node[card,fill=redbg,draw=red,text width=2.45cm] (shortcut) at (9.55,4.85) {找到捷徑};
\node[card,fill=greenbg,draw=green,text width=2.45cm] (audit) at (13.15,4.85) {外部查核};

\draw[->,draw=navy,line width=1.05pt] (signal)--(opt);
\draw[->,draw=navy,line width=1.05pt] (opt)--(shortcut);
\draw[->,draw=navy,line width=1.05pt] (shortcut)--(audit);

\node[card,fill=white,draw=line,text width=10.8cm] at (8,2.45) {只看測試、LLM 評分或 benchmark，都可能把目標縮得太窄。};
\node[note,text width=11.5cm] at (8,.95) {分數是代理，不是目的；代理一旦失真，系統越努力，偏差越穩定。};
\end{tikzpicture}
\end{document}
